\documentclass[11pt,a4paper]{article}
\usepackage[utf8]{inputenc}
\usepackage[T1]{fontenc}
\usepackage[brazil]{babel}
\usepackage[margin=2.5cm]{geometry}
\usepackage{mathptmx}
\usepackage{enumitem}
\usepackage{booktabs}
\usepackage{tabularx}
\usepackage{array}
\usepackage{hyperref}
\hypersetup{colorlinks=false, pdfborder={0 0 0}}
\usepackage{fancyhdr}

\pagestyle{fancy}
\fancyhf{}
\fancyhead[L]{\small TrackFleet --- Termo da Equipe}
\fancyhead[R]{\small\thepage}
\renewcommand{\headrulewidth}{0.4pt}

\newcolumntype{L}{>{\raggedright\arraybackslash}X}

\begin{document}

\begin{center}
{\Large\bfseries DOCUMENTO DE RECURSOS DO PROJETO}\\[2pt]
{\large\bfseries TrackFleet --- Gestão de Rastreador de Automóveis}\\[4pt]
Disciplina de Gerenciamento de Projetos $\cdot$ Framework PMBOK\textregistered\ 8\textsuperscript{a} Edição $\cdot$ Domínio: Recursos\\[2pt]
Integrantes: Enzo Allebrand e Kerlon Ribeiro (dois GPs) $\cdot$ 4\textsuperscript{o} semestre $\cdot$ Data: 03/09
\end{center}

\vspace{2mm}
\hrule
\vspace{4mm}

\noindent\textbf{\large Termo da Equipe (Team Charter)}

\vspace{1mm}
\noindent Documenta os valores, as regras de trabalho e as estratégias de comunicação e liderança da equipe, fundamentais para manter a performance e mitigar atritos em um projeto com recursos enxutos.

\section{Acordos de Trabalho e Convivência}

\begin{itemize}[itemsep=4pt]
    \item \textbf{Comunicação e Sincronização (Rituais):} 
    A comunicação será primariamente assíncrona, via Discord/WhatsApp. Será adotada uma breve sincronização ("Daily" adaptada) de 15 minutos às 19h para alinhar prioridades e apontar bloqueios.
    \item \textbf{Gestão de Qualidade (Code Review):} 
    É regra mandatória que nenhum código seja integrado (merge) na branch principal sem uma revisão do outro membro. Se Enzo codifica, Kerlon revisa e aprova, e vice-versa.
    \item \textbf{Resolução de Conflitos:} 
    A abordagem preferencial é a \textbf{Colaboração/Resolução} (busca de solução ganha-ganha, baseada em dados técnicos). Sendo uma co-gestão sem hierarquia, em caso de impasse técnico prolongado que afete o prazo, o Patrocinador (ou Orientador) atuará como fiel da balança.
\end{itemize}

\section{Fase Atual da Equipe (Modelo de Tuckman)}

Considerando a relação de pares já estabelecida em outros trabalhos acadêmicos, a equipe avança rapidamente pelas fases. No momento da assinatura deste termo, o time se encontra na transição da fase de \textbf{Normatização} para \textbf{Desempenho}:
\begin{itemize}[itemsep=2pt]
    \item Os papéis cruzados já foram debatidos e aceitos.
    \item As regras de código e responsabilidades da Matriz RACI servem para evitar a fase de \textit{Confronto} tardia.
    \item A equipe atua como auto-organizada, com foco em produtividade e entregas funcionais.
\end{itemize}

\end{document}
